\section{The fixed whole-chart square of the pushed canonical differential}

Let $\pi:Y\to X$ be a morphism of schemes and let $C$ be an actual whole
free chart of $\pi$: an affine open $U\subseteq X$ with open immersion
$g:A=\operatorname{Spec}\Gamma(X,U)\to X$, the coefficient algebra $B$ of
$C$ with scalar map $p:\operatorname{Spec}B\to A$, and the whole source
open immersion $v:\operatorname{Spec}B\to Y$ onto $\pi^{-1}(U)$. These
form the Cartesian square $g\circ p=\pi\circ v$. Fix a field $k$, a
morphism $f:X\to\operatorname{Spec}k$ and $n\geq0$. As in the previous
sections, $\omega_h^{(n)}$ is the $n$th exterior cotangent module of a
morphism $h$ to $\operatorname{Spec}k$, and for a morphism $a$ the map
$d_a:a^*\omega_h^{(n)}\to\omega_{h\circ a}^{(n)}$ is the exterior
pullback differential. For an open immersion $j$, $\operatorname{Res}_j$
is the restriction functor and $\rho_j:\operatorname{Res}_j\simeq j^*$ the
fixed restriction--pullback comparison. The Cartesian restriction--pushforward
comparison is
$\kappa:\operatorname{Res}_g\circ\pi_*\simeq p_*\circ\operatorname{Res}_v$.
Let $\chi:p^*g^*\omega_f^{(n)}\simeq v^*\pi^*\omega_f^{(n)}$ be the native
comparison formed by the pullback-composition isomorphism for $(p,g)$, the
pullback congruence for the equality $g\circ p=\pi\circ v$, and the inverse
pullback-composition isomorphism for $(v,\pi)$, and let
$H:\omega_{f\circ\pi\circ v}^{(n)}\to\omega_{f\circ g\circ p}^{(n)}$ be the
transport along the equality $f\circ\pi\circ v=f\circ g\circ p$. In the
formulas below, $u;w$ means first $u$, then $w$.

\begin{definition}[The fixed pulled canonical comparison on a whole chart]
\label{def:kn-native-chart-pulled-canonical-comparison}
\lean{EnvelopeProperty.Vanishing.NadelGlobalStructureTrace.WholeFreeChart.canonicalPullbackRestrictionIso}
\uses{sc0-open-differential:pullback-isomorphism,lem:kn-structure-trace-chart-cartesian}
\leanok
The fixed pulled canonical comparison is
\[
 \Phi_A=\kappa_{\pi^*\omega_f^{(n)}};\,
 p_*\bigl(\rho_{v,\pi^*\omega_f^{(n)}};\chi^{-1};p^*d_g\bigr):
 \operatorname{Res}_g\pi_*\pi^*\omega_f^{(n)}\simeq
 p_*p^*\omega_{f\circ g}^{(n)},
\]
where $d_g$ is invertible because $g$ is an open immersion.
\end{definition}

\begin{lemma}[The fixed chart comparisons as explicit isomorphism words]
\label{lem:kn-native-chart-comparison-words}
\lean{EnvelopeProperty.Vanishing.NadelGlobalStructureTrace.WholeFreeChart.sourceMap_comp_eq,EnvelopeProperty.Vanishing.NadelGlobalStructureTrace.WholeFreeChart.canonicalSourceRestrictionIso_eq_trans,EnvelopeProperty.Vanishing.NadelGlobalStructureTrace.WholeFreeChart.canonicalPullbackRestrictionIso_eq_trans}
\uses{def:kn-canonical-whole-chart-fixed-native-comparisons,def:kn-native-chart-pulled-canonical-comparison,lem:kn-structure-trace-chart-cartesian}
\leanok
The two composites $\operatorname{Spec}B\to\operatorname{Spec}k$ agree:
$f\circ\pi\circ v=f\circ g\circ p$. The fixed source canonical comparison
$\Phi_B:\operatorname{Res}_g\pi_*\omega_{f\circ\pi}^{(n)}\simeq
p_*\omega_{f\circ g\circ p}^{(n)}$ equals
$\kappa_{\omega_{f\circ\pi}^{(n)}};p_*(e_B)$ with
$e_B=\rho_{v,\omega_{f\circ\pi}^{(n)}};d_v;H$, and the fixed pulled canonical
comparison equals $\kappa_{\pi^*\omega_f^{(n)}};p_*(e_A)$ with
$e_A=\rho_{v,\pi^*\omega_f^{(n)}};\chi^{-1};p^*d_g$; here $d_v$, $d_g$ and $H$
are regarded as isomorphisms.
\end{lemma}
\begin{proof}
The equality of composites follows from $g\circ p=\pi\circ v$ by
associativity. The two descriptions of the comparisons are their definitions:
each is $\kappa$ at the relevant object followed by the pushforward of the
displayed composite of isomorphisms.
\end{proof}

\begin{theorem}[The fixed whole-chart differential square]
\label{thm:kn-native-chart-differential-square}
\lean{EnvelopeProperty.Vanishing.NadelGlobalStructureTrace.WholeFreeChart.canonicalSourceRestrictionIso_differential,EnvelopeProperty.Vanishing.NadelGlobalStructureTrace.WholeFreeChart.canonicalSourceRestrictionIso_differential_assoc}
\uses{lem:kn-native-chart-comparison-words,thm:kn-native-restricted-differential-general-square,lem:kn-native-restricted-differential-functorial-assembly}
\leanok
With $\Phi_A,\Phi_B$ the fixed pulled and source canonical comparisons,
\[
 \operatorname{Res}_g\bigl(\pi_*d_\pi\bigr);\,\Phi_B
 =\Phi_A;\,p_*\bigl(d_p\bigr),
\]
where on the right $d_p:p^*\omega_{f\circ g}^{(n)}\to
\omega_{f\circ g\circ p}^{(n)}$ is the chart differential. No comparison
isomorphism and no commuting square is an input. The same equality holds
after composing both sides with an arbitrary morphism out of
$p_*\omega_{f\circ g\circ p}^{(n)}$.
\end{theorem}
\begin{proof}
Write $b=\rho_{v,\omega_{f\circ\pi}^{(n)}};d_v;H$ and
$a=\rho_{v,\pi^*\omega_f^{(n)}};\chi^{-1};p^*d_g$ for the composites of
morphisms, where $\chi^{-1}$ is written as the original inverse word: the
composition comparison for $(v,\pi)$, the inverse of the congruence for
$g\circ p=\pi\circ v$, and the inverse composition comparison for $(p,g)$.
The restricted native differential square, applied to the arrows $g,\pi,p,v$
of the chart (with $v$ an open immersion), gives
\[
 \operatorname{Res}_v(d_\pi);\,b=a;\,d_p .
\]
The homomorphism of $e_B$ is $b$, since the homomorphism of a composite of
isomorphisms is the composite of their homomorphisms. For three composable
isomorphisms $r,c,t$ one has $(r;c;t^{-1})^{-1}=t;c^{-1};r^{-1}$, and the
inverses of the components of a natural isomorphism are the components of
the inverse transformation; hence the homomorphism of $e_A$ is $a$.

Apply $p_*$ to the restricted square:
$p_*\operatorname{Res}_v(d_\pi);p_*(b)=p_*(a);p_*(d_p)$. Naturality of the
Cartesian comparison $\kappa$ at $d_\pi$ gives
$\operatorname{Res}_g(\pi_*d_\pi);\kappa_{\omega_{f\circ\pi}^{(n)}}
=\kappa_{\pi^*\omega_f^{(n)}};p_*\operatorname{Res}_v(d_\pi)$. Combining the
two equalities,
\[
 \operatorname{Res}_g(\pi_*d_\pi);\kappa_{\omega_{f\circ\pi}^{(n)}};p_*(b)
 =\kappa_{\pi^*\omega_f^{(n)}};p_*(a);p_*(d_p).
\]
By the preceding lemma and the identification of the homomorphisms of
$e_B,e_A$ with $b,a$, the left factor
$\kappa_{\omega_{f\circ\pi}^{(n)}};p_*(b)$ is the homomorphism of $\Phi_B$ and
$\kappa_{\pi^*\omega_f^{(n)}};p_*(a)$ is the homomorphism of $\Phi_A$. This is
the asserted equality; it is the functorial assembly lemma applied with
$F_1=\pi_*$, $F_2=\operatorname{Res}_g$, $G_1=\operatorname{Res}_v$,
$G_2=p_*$ and $e=\kappa$. Composing both sides on the right with a further
morphism and reassociating gives the whiskered form.
\end{proof}
